\documentclass[11pt]{article}
\usepackage[a4paper,margin=1in]{geometry}
\usepackage{amsmath,amssymb}
\usepackage{hyperref}
\title{Orientation Representations and Plane--Direction Construction in PyTex}
\author{PyTex Contributors}
\date{\today}

\begin{document}
\maketitle

\section{Scope}

This note records two related PyTex foundations:

\begin{itemize}
  \item vectorized conversion among rotation matrices, unit quaternions, axis--angle pairs, and Rodrigues / Rodrigues--Frank coordinates,
  \item construction of crystal-to-specimen orientations from a crystal plane and a crystal direction.
\end{itemize}

The implementation is intentionally explicit about mapping direction, singularities, and default specimen references so the batch APIs remain scientifically auditable.

\section{Canonical Rotation Representation Rules}

PyTex stores unit quaternions in $(w, x, y, z)$ order and interprets them as active crystal-to-specimen rotations. For an axis--angle pair $(\hat{\mathbf{n}}, \omega)$ the quaternion is

\[
q = \left(\cos\frac{\omega}{2},\ \hat{\mathbf{n}}\sin\frac{\omega}{2}\right).
\]

For batch export to axis--angle form, PyTex canonicalizes the quaternion sign so $w \ge 0$ before extracting the angle. This yields

\[
\omega \in [0, \pi],
\]

with a default axis of $\hat{\mathbf{e}}_z$ for the identity rotation. That default axis is a storage convention only; the identity rotation has no unique physical axis.

\section{Rodrigues and Rodrigues--Frank Coordinates}

PyTex uses the Rowenhorst et al.\ definition of Rodrigues coordinates:

\[
\boldsymbol{\rho} = \hat{\mathbf{n}}\tan\frac{\omega}{2}.
\]

The Rodrigues--Frank form stores the same information as a four-component vector

\[
\left(\hat{\mathbf{n}}, \tan\frac{\omega}{2}\right).
\]

This makes the two-fold singularity explicit. When $\omega = \pi$, the Rodrigues vector magnitude tends to infinity and the Rodrigues--Frank scale is represented as $+\infty$ in PyTex. That representation is deliberate: it preserves the mathematical singularity instead of hiding it behind an arbitrary clip value.

\section{Plane--Direction Orientation Construction}

Given a crystal plane normal $\hat{\mathbf{n}}_c$ and a crystal direction $\mathbf{d}_c$, PyTex first projects the direction into the plane:

\[
\hat{\mathbf{x}}_c =
\frac{\mathbf{d}_c - (\mathbf{d}_c \cdot \hat{\mathbf{n}}_c)\hat{\mathbf{n}}_c}
{\left\lVert \mathbf{d}_c - (\mathbf{d}_c \cdot \hat{\mathbf{n}}_c)\hat{\mathbf{n}}_c \right\rVert}.
\]

The corresponding right-handed crystal basis is

\[
\hat{\mathbf{y}}_c = \hat{\mathbf{n}}_c \times \hat{\mathbf{x}}_c, \qquad
\hat{\mathbf{z}}_c = \hat{\mathbf{n}}_c.
\]

PyTex constructs an analogous specimen basis from the requested specimen plane normal and specimen in-plane direction. The current default is:

\begin{itemize}
  \item crystal plane normal $\rightarrow$ specimen $Z$,
  \item in-plane crystal direction $\rightarrow$ specimen $X$.
\end{itemize}

The rotation matrix is then

\[
\mathbf{R}_{c \rightarrow s} =
\mathbf{B}_s \mathbf{B}_c^{\mathsf{T}},
\]

where $\mathbf{B}_c = [\hat{\mathbf{x}}_c\ \hat{\mathbf{y}}_c\ \hat{\mathbf{z}}_c]$ and
$\mathbf{B}_s = [\hat{\mathbf{x}}_s\ \hat{\mathbf{y}}_s\ \hat{\mathbf{z}}_s]$.

This guarantees

\[
\mathbf{R}_{c \rightarrow s}\hat{\mathbf{n}}_c = \hat{\mathbf{n}}_s,
\qquad
\mathbf{R}_{c \rightarrow s}\hat{\mathbf{x}}_c = \hat{\mathbf{x}}_s,
\]

and preserves the PyTex orientation convention that orientations map crystal vectors into specimen vectors.

\section{Failure Modes}

PyTex rejects the construction if:

\begin{itemize}
  \item the supplied plane and direction do not share a phase,
  \item the projected crystal direction becomes zero,
  \item the projected specimen direction becomes zero,
  \item the resulting triads do not define proper right-handed bases.
\end{itemize}

These checks are performed at construction time so invalid frame or phase semantics cannot survive into later EBSD, PF, IPF, or ODF workflows.

\section{Normative References}

\begin{thebibliography}{9}
\bibitem{bunge}
H.-J. Bunge,
\textit{Texture Analysis in Materials Science: Mathematical Methods},
Butterworths, 1969.
DOI: \url{https://doi.org/10.1016/C2013-0-11769-2}.

\bibitem{rowenhorst}
D. Rowenhorst et al.,
``Consistent Representations of and Conversions Between 3D Rotations'',
\textit{Modelling and Simulation in Materials Science and Engineering}, 23(8), 2015.
DOI: \url{https://doi.org/10.1088/0965-0393/23/8/083501}.
\end{thebibliography}

\end{document}
